% =====================================================================
%  config/estilo.tex — Diseño y entornos propios del manual
%  Requiere config/paquetes.tex cargado antes.
% =====================================================================
\makeatletter

% ---------------------------------------------------------------------
% 1. Página
% ---------------------------------------------------------------------
\geometry{
  paperwidth  = 170mm,
  paperheight = 240mm,
  inner       = 20mm,
  outer       = 22mm,
  top         = 24mm,
  bottom      = 26mm,
  headsep     = 7mm,
  footskip    = 11mm,
}
\setlength{\parindent}{1.2em}
\linespread{1.06}
\frenchspacing
\setlength{\emergencystretch}{1.5em}
\raggedbottom   % sin estirar el espacio vertical entre ejemplos y recuadros

% ---------------------------------------------------------------------
% 2. Fuentes (Libertinus: incluida en TeX Live, con versalitas reales)
% ---------------------------------------------------------------------
\setmainfont{LibertinusSerif}[
  Extension      = .otf,
  UprightFont    = *-Regular,
  ItalicFont     = *-Italic,
  BoldFont       = *-Semibold,
  BoldItalicFont = *-SemiboldItalic,
  Ligatures      = TeX,
  Numbers        = Lining]
\setsansfont{LibertinusSans}[
  Extension      = .otf,
  UprightFont    = *-Regular,
  ItalicFont     = *-Italic,
  BoldFont       = *-Bold,
  Ligatures      = TeX,
  Scale          = MatchLowercase]
\newfontfamily\versalitasespaciadas{LibertinusSans}[
  Extension = .otf, UprightFont = *-Regular, BoldFont = *-Bold,
  LetterSpace = 12, Scale = MatchLowercase]

% ---------------------------------------------------------------------
% 3. Colores
% ---------------------------------------------------------------------
\definecolor{acento}{HTML}{1B5E6B}       % petróleo: títulos, ejercicios
\definecolor{acentoclaro}{HTML}{EAF2F3}  % fondo de recuadros
\definecolor{acento2}{HTML}{A4502A}      % terracota: ejemplos resueltos
\definecolor{gristexto}{HTML}{555555}

% ---------------------------------------------------------------------
% 4. Títulos (KOMA-Script)
% ---------------------------------------------------------------------
\setkomafont{disposition}{\sffamily\bfseries\color{acento}}
\setkomafont{chapterprefix}{\normalsize\versalitasespaciadas\bfseries\color{acento}}
\setkomafont{chapter}{\sffamily\bfseries\LARGE\color{black}}
\setkomafont{section}{\sffamily\bfseries\Large}
\setkomafont{subsection}{\sffamily\bfseries\large}
\setkomafont{subsubsection}{\sffamily\bfseries\normalsize}

\renewcommand*{\chapterformat}{\MakeUppercase{\chaptername}~\thechapter}
\renewcommand*{\chapterlinesformat}[3]{%
  \Ifstr{#1}{chapter}{%
    \raggedright
    \Ifstr{#2}{}{}{{\usekomafont{chapterprefix}#2}\par\nobreak\vspace{.5\baselineskip}}%
    #3\par
  }{\@hangfrom{#2}{#3}}}
\RedeclareSectionCommand[beforeskip=10mm, afterskip=4mm]{chapter}
\RedeclareSectionCommand[beforeskip=-2.2ex plus -.5ex minus -.2ex,
                         afterskip=1.1ex plus .2ex]{section}
\RedeclareSectionCommand[beforeskip=-1.8ex plus -.4ex minus -.2ex,
                         afterskip=.8ex plus .1ex]{subsection}
\setcounter{secnumdepth}{\subsubsectionnumdepth}
% Secciones numeradas con el número de capítulo: 2.1, 2.3.1
\setcounter{tocdepth}{\sectiontocdepth}

% \capitulo[título corto]{Título}{Autor/es}
%   Abre un capítulo numerado con firma de autor (en el cuerpo y en el índice).
\NewDocumentCommand{\capitulo}{o m m}{%
  \IfNoValueTF{#1}{\chapter{#2}}{\chapter[#1]{#2}}%
  \label{cap:\capclave}%
  \xappto\capsincluidos{\capclave,}%
  \gdef\autorcapitulo{#3}%
  {\raggedright\normalfont\large\itshape #3\par}%
  \vspace{2mm}\noindent{\color{acento}\rule{\linewidth}{.6pt}}\par
  \vspace{7mm}%
  \addtocontents{toc}{\protect\autorentoc{#3}}%
  \@afterindentfalse\@afterheading}
\newcommand*{\autorcapitulo}{}
\newcommand*{\autorentoc}[1]{%
  \par\nopagebreak
  {\leftskip=1.5em\relax\rightskip=\@tocrmarg\relax
   \small\itshape\color{gristexto}#1\par}%
  \addvspace{2pt}}

% Secciones finales de capítulo (no numeradas, fuera del índice,
% con marcador PDF): Lecturas, Solucionario, Glosario, Referencias.
\newcounter{seccionfinal}
\newcommand*{\seccionfinal}[1]{%
  \stepcounter{seccionfinal}%
  \addsec*{#1}%
  \pdfbookmark[1]{#1}{seccionfinal.\theseccionfinal}%
  \markright{#1}}

% ---------------------------------------------------------------------
% 5. Encabezados y pies
% ---------------------------------------------------------------------
\setkomafont{pageheadfoot}{\small\sffamily\color{gristexto}}
\setkomafont{pagenumber}{\small\sffamily\bfseries\color{acento}}
\renewcommand*{\chaptermarkformat}{\thechapter.\enskip}
\renewcommand*{\sectionmarkformat}{\thesection\enskip}
\automark[section]{chapter}
\clearpairofpagestyles
\ohead{\pagemark}
\ihead{\headmark}
\cfoot[\pagemark]{}   % páginas «plain» (apertura de capítulo)
\KOMAoptions{headsepline=false}

% ---------------------------------------------------------------------
% 6. Notas al pie, epígrafes de tablas y figuras
% ---------------------------------------------------------------------
\deffootnote[1.2em]{1.2em}{1em}{\textsuperscript{\thefootnotemark}\,}
\setfootnoterule[.4pt]{3cm}
\addtokomafont{footnote}{\linespread{1}\selectfont}

\renewcommand*{\thetable}{\arabic{table}}     % «Tabla 1», reinicia por capítulo
\renewcommand*{\thefigure}{\arabic{figure}}
\renewcommand*{\captionformat}{.\enskip}
\setkomafont{caption}{\small}
\setkomafont{captionlabel}{\small\sffamily\bfseries\color{acento}}
\setcapindent{0pt}
\KOMAoptions{captions=tablesignature}         % epígrafe debajo de la tabla

% Tablas: columnas centradas y cuerpo en \small
\newcolumntype{C}{>{\centering\arraybackslash}X}
\newcolumntype{L}{>{\raggedright\arraybackslash}X}
\newcolumntype{Y}[1]{>{\centering\arraybackslash}p{#1}}
\AtBeginEnvironment{tabularx}{\small\renewcommand{\arraystretch}{1.2}}
\AtBeginEnvironment{tabular}{\small\renewcommand{\arraystretch}{1.2}}
% Valores de los cuadros de ejercicios
\newcommand*{\mas}{+}
\newcommand*{\menos}{\textendash}
\newcommand*{\masmenos}{+/\textendash}
% Encabezado de columna en cursiva (como en los originales)
\newcommand*{\enc}[1]{\textit{#1}}

% ---------------------------------------------------------------------
% 7. Ejemplos (gb4e)
%    Texto:      numeración continua por capítulo, entre paréntesis: (1), (2a)
%    Ejercicios: numeración propia, reinicia en cada ejercicio, sin
%                paréntesis: 1., 2. a.
% ---------------------------------------------------------------------
\counterwithin*{exx}{chapter}
\newcounter{ejx}
\newcommand*{\exnumcounter}{exx}
\newcommand*{\exlabelformat}[1]{(#1)}

\patchcmd{\@exe}{(\thexnumi)}{\exlabelformat{\thexnumi}}{}%
  {\PackageError{estilo}{No se pudo adaptar gb4e (\string\@exe)}{}}
\patchcmd{\@exe}{\setcounter{xnumi}{\value{exx}}}%
                {\setcounter{xnumi}{\value{\exnumcounter}}}{}%
  {\PackageError{estilo}{No se pudo adaptar gb4e (contador)}{}}
\patchcmd{\endexe}{\setcounter{exx}{\value{xnumi}}}%
                  {\setcounter{\exnumcounter}{\value{xnumi}}}{}%
  {\PackageError{estilo}{No se pudo adaptar gb4e (\string\endexe)}{}}

% Dentro de ejercicios, cada bloque exe reinicia la numeración
\pretocmd{\exe}{\ifdefstring{\exnumcounter}{ejx}{\setcounter{ejx}{0}}{}}{}%
  {\PackageError{estilo}{No se pudo adaptar gb4e (\string\exe)}{}}

% Modo «ejercicio»: contador ejx, etiqueta sin paréntesis
\newcommand*{\modoejercicio}{%
  \setcounter{ejx}{0}%
  \renewcommand*{\exnumcounter}{ejx}%
  \renewcommand*{\exlabelformat}[1]{##1.}%
  \exewidth{99.}}

% Remisiones: \exref{ex:ambig} → (7);  \exrefs{ex:a}{ex:b} → (7)-(9)
\newcommand*{\exref}[1]{(\ref{#1})}
\newcommand*{\exrefs}[2]{(\ref{#1})-(\ref{#2})}

% Juicio en línea: \juicio{*}{haiga} → *haiga
\newcommand*{\juicio}[2]{#1#2}

% Abreviaturas de glosa propias (leipzig): \Dom
\newleipzig{dom}{dom}{marcado diferencial de objeto}

% Formas mencionadas y significados
\newcommand*{\forma}[1]{\textit{#1}}               % lengua objeto
\newcommand*{\sig}[1]{\textquoteleft#1\textquoteright} % ‘significado’
\newcommand*{\termino}[1]{\textit{#1}}             % término técnico introducido
\newcommand*{\concepto}[1]{\textsc{\MakeLowercase{#1}}} % VERSALITA de las citas

% Bloques de cita (APA, ≥ 40 palabras)
\newenvironment{cita}
  {\begin{quote}\small\linespread{1}\selectfont}
  {\end{quote}}

% ---------------------------------------------------------------------
% 8. Árboles (forest)
% ---------------------------------------------------------------------
\forestset{
  gramatica/.style={
    for tree={
      parent anchor=south, child anchor=north,
      s sep=6mm, l sep=4mm, inner sep=1pt,
      font=\small, align=center,
    },
    baseline,   % alinea la raíz con la etiqueta del ejemplo
  },
  triangulo/.style={roof},
}
\NewDocumentEnvironment{arbol}{+b}
  {\par\medskip\begin{center}\begin{forest}gramatica #1\end{forest}\end{center}\medskip}{}

% ---------------------------------------------------------------------
% 9. Ejercicios
%    \begin{ejercicio}[subtítulo opcional] ... \end{ejercicio}
%    \begin{resuelto}[Ejemplo resuelto] ... \end{resuelto}
% ---------------------------------------------------------------------
\newcounter{ejercicio}[chapter]
\renewcommand*{\theejercicio}{\arabic{ejercicio}}

\newtcolorbox{cajaejercicio}[1]{%
  enhanced, breakable, blanker,
  borderline west={2.2pt}{0pt}{acento},
  left=4mm, right=0mm, top=0mm, bottom=1mm,
  before skip=6mm plus 1mm, after skip=6mm plus 1mm,
  fontupper=\small,
  title={#1}, coltitle=acento,
  fonttitle=\sffamily\bfseries,
  attach title to upper={\par\smallskip},
  parbox=false,
}
\NewDocumentEnvironment{ejercicio}{o}
  {\par\needspace{8\baselineskip}\refstepcounter{ejercicio}%
   \begin{cajaejercicio}{Ejercicio~\theejercicio%
       \IfValueT{#1}{\enskip{\mdseries\color{gristexto}·\enskip #1}}}%
   \modoejercicio\setlength{\parindent}{0pt}\setlength{\parskip}{.5\baselineskip}}
  {\end{cajaejercicio}}

\NewDocumentEnvironment{resuelto}{O{Ejemplo resuelto}}
  {\par\smallskip\small\noindent
   {\sffamily\bfseries\color{acento2}#1.}\enskip\ignorespaces}
  {\par}

% Listas del texto: espaciado compacto
\setlist[itemize]{itemsep=.25\baselineskip, topsep=.4\baselineskip,
                  parsep=0pt, leftmargin=1.6em}
\setlist[itemize,2]{label=--, topsep=.2\baselineskip}
\setlist[enumerate]{itemsep=.25\baselineskip, topsep=.4\baselineskip, parsep=0pt}

% Preguntas o consignas internas del ejercicio: a), b)...
\newlist{consignas}{enumerate}{1}
\setlist[consignas]{label=\alph*), leftmargin=1.8em, topsep=.3\baselineskip,
                    itemsep=.2\baselineskip}

% ---------------------------------------------------------------------
% 10. Recuadros destacados
%    \begin{recuadro}[Título] ... \end{recuadro}  (con título)
%    \begin{destacado} ... \end{destacado}         (sin título)
% ---------------------------------------------------------------------
\newtcolorbox{recuadro}[1][]{%
  enhanced, breakable,
  colback=acentoclaro, colframe=acento,
  boxrule=0pt, toprule=1.4pt, bottomrule=0pt, leftrule=0pt, rightrule=0pt,
  arc=0pt, outer arc=0pt,
  left=5mm, right=5mm, top=3mm, bottom=3mm,
  before skip=5mm, after skip=5mm,
  fontupper=\small,
  title={#1}, coltitle=acento, fonttitle=\sffamily\bfseries,
  attach title to upper={\par\smallskip},
}
\newtcolorbox{destacado}{%
  enhanced, breakable, colback=acentoclaro, frame hidden,
  arc=0pt, left=5mm, right=5mm, top=2.5mm, bottom=2.5mm,
  before skip=4mm, after skip=4mm,
}

% ---------------------------------------------------------------------
% 11. Secciones finales del capítulo
% ---------------------------------------------------------------------
% Clave del capítulo para filtrar la bibliografía: cap01, cap02...
\newcommand*{\capclave}{\two@digits{\value{chapter}}}

\defbibheading{seccionfinal}[]{\seccionfinal{#1}}

% Bibliografía general (al final del libro): entradas de los capítulos
% incluidos (palabras clave capNN y lectNN) más todo lo citado con \cite.
% Las entradas marcadas «nobib» (duplicados para lecturas) se omiten.
\newcommand*{\capsincluidos}{}
\DeclareBibliographyCategory{citados}
\AtEveryCitekey{\addtocategory{citados}{\thefield{entrykey}}}
\newcommand*{\bib@chequearcap}[1]{%
  \ifkeyword{cap#1}{\def\bib@ok{1}}{}%
  \ifkeyword{lect#1}{\def\bib@ok{1}}{}}
\defbibcheck{libro}{%
  \ifkeyword{nobib}
    {\skipentry}
    {\def\bib@ok{0}%
     \ifcategory{citados}{\def\bib@ok{1}}{}%
     \forcsvlist{\bib@chequearcap}{\capsincluidos}%
     \ifnum\bib@ok=0 \skipentry\fi}}
\newcommand*{\bibliografiageneral}{%
  \printbibliography[check=libro, heading=bibintoc,
                     title={Referencias bibliográficas}]}

\newcommand*{\lecturasrecomendadas}{%
  \begingroup
  \edef\x{\noexpand\printbibliography[keyword=lect\capclave,
          heading=seccionfinal, title={Lecturas recomendadas}]}%
  \x\endgroup}

% Tipos de editor en español: editortype = {coordinator|director|compiler}
%   → (Coord./Coords.), (Dir./Dirs.), (Comp./Comps.)
\csdef{tipoeditor@coordinator}{Coord\adddot}
\csdef{tipoeditor@coordinators}{Coords\adddot}
\csdef{tipoeditor@director}{Dir\adddot}
\csdef{tipoeditor@directors}{Dirs\adddot}
\csdef{tipoeditor@compiler}{Comp\adddot}
\csdef{tipoeditor@compilers}{Comps\adddot}
\renewbibmacro*{apaeditorstrg}[1]{%
  \iffieldundef{#1type}
    {\ifthenelse{\value{#1}>1\OR\ifandothers{#1}}
       {\bibcpstring{editors}}
       {\bibcpstring{editor}}}
    {\edef\tipoed@{\thefield{#1type}}%
     \ifthenelse{\value{#1}>1\OR\ifandothers{#1}}
       {\ifcsdef{tipoeditor@\tipoed@ s}
          {\csuse{tipoeditor@\tipoed@ s}}
          {\bibcpstring{type\thefield{#1type}s}}}
       {\ifcsdef{tipoeditor@\tipoed@}
          {\csuse{tipoeditor@\tipoed@}}
          {\bibcpstring{type\thefield{#1type}}}}}}

% APA en español: «y» en lugar de «&»; «Obra original publicada en»;
% partícula pospuesta («Saussure, F. de»)
\ExecuteBibliographyOptions{useprefix=false}
\DeclareDelimFormat[bib,biblist]{finalnamedelim}{%
  \ifthenelse{\value{listcount}>\maxprtauth}
    {}
    {\addspace y\space}}
\DeclareDelimFormat[parencite,fullcite,fullcitebib]{finalnamedelim}{\addspace y\space}
\renewbibmacro*{origyear}{%
  \ifboolexpr{%
    not test {\iffieldundef{origyear}}
    and not test {\iffieldundef{labelyear}}
    and not test {\iffieldsequal{labelyear}{origyear}}}
    {\printtext[parens]{Obra original publicada en\addnbspace\printorigdate}%
     \renewcommand*{\finentrypunct}{\relax}}
    {}}

% Dos autores sin coma antes de «y» (Bosque, I. y Demonte, V.); con tres
% o más, coma entre todos salvo antes de «y».
\renewcommand*{\apablx@ifrevnameappcomma}{%
  \iftoggle{apablx@wantcomma}
    {\ifboolexpr{test {\ifnumless{\value{listcount}}{\value{liststop}-1}}
                 or test \ifmorenames}}
    {\@secondoftwo}}

% Capítulos de libro: (Vol. 2, pp. 23-45), como en APA 7
\renewbibmacro*{addinfo1}{%
  \printtexte[parens]{%
    \usebibmacro{editor+trans+narr}%
    \setunit*{\addsemicolon\addspace}%
    \printfield{edition}%
    \setunit*{\addcomma\addspace}%
    \ifthenelse{\ifentrytype{report}\OR\ifentrytype{unpublished}\OR
                \ifentrytype{dataset}\OR\ifentrytype{software}\OR
                \ifentrytype{audio}\OR\ifentrytype{patent}\OR
                \ifentrytype{image}}{\usebibmacro{addinfo3}}{}%
    \setunit*{\addcomma\addspace}%
    \ifthenelse{\ifentrytype{book}\OR\ifentrytype{inbook}}{\printfield{volumes}}{}%
    \setunit*{\addcomma\addspace}%
    \ifthenelse{\ifentrytype{book}\OR\ifentrytype{inbook}}{\printfield{volume}}{}%
    \setunit*{\addcomma\addspace}%
    \ifentrytype{inbook}{\printfield{pages}}{}}}

% Bibliografía: sangría francesa y cuerpo menor
\setlength{\bibhang}{1.5em}
\renewcommand*{\bibfont}{\small\raggedright}

% Solucionario
\newcommand*{\solucionario}{\seccionfinal{Solucionario}}
\NewDocumentEnvironment{solucion}{m}
  {\par\addvspace{\medskipamount}%
   \needspace{4\baselineskip}%
   \noindent{\sffamily\bfseries\color{acento}Ejercicio~#1}\par\nobreak
   \small\modoejercicio\setlength{\parindent}{0pt}\setlength{\parskip}{.4\baselineskip}}
  {\par}

% Glosario (por capítulo y general)
\newlist{glosario}{description}{4}
\setlist[glosario]{font=\normalfont\bfseries\itshape, leftmargin=1.5em,
                   labelindent=0pt, itemsep=.35\baselineskip,
                   style=sameline, before=\small}
\setlist[glosario,2]{leftmargin=1.2em, labelindent=0pt, topsep=.2\baselineskip}
\setlist[glosario,3]{leftmargin=1.2em, labelindent=0pt, topsep=.2\baselineskip}
\setlist[glosario,4]{leftmargin=1.2em, labelindent=0pt, topsep=.2\baselineskip}
% Término con subentradas: \item[Distribución]\subterminos  (fuerza salto de línea)
\newcommand*{\subterminos}{\mbox{}}
\newcommand*{\glosariocapitulo}{\seccionfinal{Glosario}}

% Glosario general (finales/glosario.tex): dos columnas
%   \begin{glosariogeneral} \item[Término:] Definición. \end{glosariogeneral}
\setlength{\columnsep}{7mm}
\setlength{\columnseprule}{0pt}
\NewDocumentEnvironment{glosariogeneral}{}
  {\begin{multicols}{2}\raggedcolumns\begin{glosario}[before=\small\raggedright]}
  {\end{glosario}\end{multicols}}

% Siglario
\newlist{siglario}{description}{1}
\setlist[siglario]{font=\normalfont\sffamily\bfseries, labelwidth=4.5em,
                   leftmargin=5.5em, labelsep=1em, align=left,
                   itemsep=.15\baselineskip, style=standard}

% ---------------------------------------------------------------------
% 12. Índice general
% ---------------------------------------------------------------------
\setkomafont{chapterentry}{\sffamily\bfseries}
\RedeclareSectionCommand[tocindent=0pt, tocnumwidth=1.5em]{chapter}
\RedeclareSectionCommand[tocindent=1.5em, tocnumwidth=2.3em]{section}


% ---------------------------------------------------------------------
% 13. Utilidades de edición y de ejercicios
% ---------------------------------------------------------------------
% Marca visible de algo a revisar en el original (buscar «\pendiente»)
\newcommand*{\pendiente}[1]{{\color{red!70!black}\sffamily\small[#1]}}

% Respuesta de un ítem resuelto dentro de un ejercicio
\newcommand*{\respuesta}[1]{{\color{acento2}\itshape #1}}

% Líneas para completar (número de renglones, por defecto 2)
\NewDocumentCommand{\completar}{O{2}}{%
  \enspace\dotfill\mbox{}%
  \foreach \i in {2,...,#1}{\newline\mbox{}\dotfill\mbox{}}}

% Palabra subrayada en consignas
\newcommand*{\subrayado}[1]{\uline{#1}}

% Texto de lectura (cuentos, fragmentos para analizar)
\newtcolorbox{textolectura}[1][]{%
  enhanced, breakable,
  colback=white, colframe=acento!60, boxrule=.5pt, arc=0pt,
  left=6mm, right=6mm, top=4mm, bottom=4mm,
  before skip=5mm, after skip=5mm,
  parbox=false,
  before upper={{\noindent\sffamily\bfseries\small\color{acento}#1\par}%
                \medskip\setlength{\parindent}{1.2em}},
}

% Figura cuya imagen todavía no está en figuras/: deja un recuadro con
% el epígrafe definitivo. Uso: \figurapendiente{Epígrafe}{fig:etiqueta}
\NewDocumentCommand{\figurapendiente}{m m O{45mm}}{%
  \begin{figure}[htbp]
    \centering
    \fcolorbox{gristexto}{white}{\parbox[c][#3][c]{.8\linewidth}{\centering
      \sffamily\small\color{gristexto}Imagen pendiente\\[2pt]\footnotesize #1}}
    \caption{#1}\label{#2}
  \end{figure}}

% ---------------------------------------------------------------------
% 14. Enlaces
% ---------------------------------------------------------------------
\hypersetup{
  linkcolor = acento,     % ejemplos, secciones, notas, índice
  citecolor = acento2,    % citas bibliográficas
  urlcolor  = acento,     % direcciones web
}

% Remisión a otro capítulo: \capref{4} → «capítulo 4», con enlace si ese
% capítulo ya está incluido en el libro; texto plano si todavía no.
\newcommand*{\capref}[1]{%
  \@ifundefined{r@cap:\two@digits{#1}}%
    {capítulo~#1}%
    {\hyperref[cap:\two@digits{#1}]{capítulo~#1}}}
% Variante abreviada: \Capref{4} → «cap. 4»
\newcommand*{\Capref}[1]{%
  \@ifundefined{r@cap:\two@digits{#1}}%
    {cap.~#1}%
    {\hyperref[cap:\two@digits{#1}]{cap.~#1}}}

% Etiquetas de corchetes y nudos con barra:
%   [\etq{SN} el\etq{D/Esp} libro\etq{N/n}]   N\prima (N')
\newcommand*{\etq}[1]{\textsubscript{#1}}
\newcommand*{\prima}{\ensuremath{'}}

% Transcripción fonética [ ] y fonológica / /
\newcommand*{\fon}[1]{[#1]}
\newcommand*{\fonol}[1]{/#1/}

% Imagen o audio que todavía no está en figuras/ (no flotante, para
% ejercicios): \imagenpendiente[alto]{Descripción}
\NewDocumentCommand{\imagenpendiente}{O{22mm} m}{%
  \par\smallskip\noindent\fcolorbox{gristexto}{white}{\parbox[c][#1][c]{\dimexpr\linewidth-2\fboxsep-2\fboxrule}{\centering
    \sffamily\footnotesize\color{gristexto}#2}}\par\smallskip}

% Texto ajeno que se incorpora desde el original (letras de canciones,
% textos literarios completos): deja un recuadro con la indicación.
\newcommand*{\textoaincorporar}[1]{%
  \imagenpendiente[18mm]{Incorporar aquí #1 desde el original}}

\makeatother
